\section{Dual-Chain SDR Prototype and Methodology}\label{sec:methodology}

\subsection{Single-TX / Dual-Independent-RX Testbed}\label{sec:testbed}

A single Zynq-7020/AD9363 SDR transmitter runs an \texttt{app\_FreeRTOS} SDK-integrated build at nominal \SI{0}{dBm} (\texttt{tx\_attenuation=10000\,mdB}). Manifests freeze \proto{} firmware build \texttt{20261119-uc1-r1} across P2 (2701--2718), P3 run 2719, and all 36 P4 runs (2731--2766); as a provenance exception, P3 pilot runs 2720--2730 logged the literal string \texttt{export FW\_BUILD}, so their firmware version cannot be confirmed from manifests alone. Because DMA completion interrupts are unreliable on this board, the transmitter actively polls DMA status with bounded timeouts and exports per-burst logs plus \texttt{UC1\_STATS final=1} over \SI{115200}{baud} UART (early P2 run 2701 omitted \texttt{final=1} and was audited via its end-of-run context). DMA completion is a submission proxy, not calibrated RF airtime or energy.

Two independent HackRF One receivers sample at \SI{4}{Msps} (AMP off, fixed serial-to-antenna bindings) with phase-frozen gains: LNA~\SI{16}{dB} and VGA~\SI{10}{dB} in P2 \texttt{R0} feeder runs (2701--2718), and LNA~\SI{16}{dB} and VGA~\SI{16}{dB} in P3/P4 \texttt{R1} wireless runs (2719--2766). Isolated host pipelines run \texttt{exadv\_rx} (\ble{} ch3 at \SI{2410}{MHz} for \modejc{}/\modesc{}; ch39 at \SI{2480}{MHz} for \modelc{}) and standard O-QPSK \texttt{std\_zigbee/zigbee\_rx.py} (\stdfour{} ch12, \SI{2410}{MHz}). Each chain independently validates its outer framing, \proto{} header, and record \texttt{app\_crc16}; for fragmented \modelc{} records, the \ble{} chain first reassembles all \fragproto{} fragments from the same attempt. Deterministic bytes seeded by $(\texttt{run\_id}, \texttt{command\_id}, \texttt{side})$ verify decoded integrity only and never reconstruct missing opposite-side bytes; cross-side matching joins validated records solely by $(\texttt{run\_id}, \texttt{command\_id})$.

\subsection{Pre-Registered Matrix and Audit Metrics}\label{sec:metrics}

Experiments proceed through three frozen phases. In \textbf{P2} (\texttt{R0} direct feeder, runs 2701--2718), we verify all 10 layout/mode branches in Table~\ref{tab:budget} ($10$ tasks/run), run shared-IQ dual decoding on recorded \ble{} IQ from Points A (2704), B (2718), and C (2707), execute four equal-budget bridge runs (2711--2714, retaining manifest-reconciliation provenance), test three fixed-secondary \texttt{time}-profile runs (2715--2717, $300$ slots/run; manual feeder disconnection on \ble{} in 2716 and \stdfour{} in 2717 logged by approximate packet-index range, with state durations measured on the disrupted receiver's local IQ index), and run deterministic offline fault-injection replay (\texttt{uc1\_h6\_replay.py}) on recorded streams from runs 2714 and 2715 (synthetic offline replay, not over-the-air anomalies). In \textbf{P3} (runs 2719--2730), two whip-antenna distances at Point~B ($K{=}1$, two reversed blocks of $30$ tasks/mode, $90$ tasks/block) are screened under pre-registered rules requiring $18\text{--}76/90$ pooled joint successes per block, \emph{at least one identical mode} at $1\text{--}29/30$ in both blocks, and no queue overflows, no capture truncation, and no significant clipping (clip-ratio saturation rule). Candidate~1 (\SI{0.20}{m}) failed Block~2 ($13/90$), selecting \SI{0.10}{m} (\texttt{R1\_candidate2\_0p10m}; $31/90$ and $24/90$, zero defects). In \textbf{P4} (runs 2731--2766), Points A $(0,0)$ and B $(26,0)$, policies $K{=}1$ and $Q{=}\SI{6000}{\micro\second}$, and modes $\{\modejc{},\modesc{},\modelc{}\}$ run across $3$ counterbalanced blocks ($100$ tasks/run, \num{3600} planned tasks; raw IQ archived on pre-designated runs 2740 and 2742).

Four audit rules govern all metrics. First, each ratio uses its pre-registered denominator (zero denominator $\to$ \texttt{N/A}): (1) \emph{within-planned-attempt joint resolution rate} uses planned logical tasks---a task counts once as \texttt{both} when \ble{} and \stdfour{} independently validate it within $K$ attempts even in different slots (e.g., $7$ of $93/300$ in Point~B $Q{=}\SI{6000}{\micro\second}$ \modelc{}), whereas \modelc{}'s \ble{} chain must collect all \fragproto{} fragments within the \emph{same} attempt; remaining tasks enter \texttt{ble\_only}, \texttt{zigbee\_only}, and \texttt{neither}, alongside first-attempt rates; (2) \emph{held-out probe coverage} uses planned probes ($24$/run); and (3) \emph{in-schedule recovery} uses in-schedule loss events on that side (\texttt{N/A (0/0)} when zero). Second, without externally calibrated RF clocks, P4 deadline status is audited separately: \num{2730}/\num{3600} jointly resolved tasks are \texttt{unknown} and \num{870}/\num{3600} are \texttt{not\_decoded}. Third, unit cost divides total submitted waveform proxy (including retries, updates, and incomplete fragments) by jointly resolved tasks. Fourth, paired comparisons report all three block differences (pp), mean, and 95\% $t$-intervals ($\text{df}{=}2$, $t_{0.975,2}{=}4.303$), distinguishing protocol-identifiable co-channel pairs (\modejc{} vs.\ \modesc{}) from channel-confounded configuration pairs (\modejc{} ch3 vs.\ \modelc{} ch39).
